% ----------------------------------------------------------------------
%  2.1.6. Экспериментальные методы
% ----------------------------------------------------------------------
\subsection{Экспериментальные методы наблюдения и исследования точечных дефектов}\label{ssec:exp-methods}

\begin{orientir}
Как измерить концентрацию дефектов, энергии их образования и миграции,
если сами дефекты атомного размера? Есть два пути: увидеть дефект
напрямую или измерить свойство кристалла, которое чувствительно к
дефектам. Большая часть сведений получена вторым, косвенным путём. Здесь
разобраны два классических метода~— сравнение изменения длины образца с
изменением параметра решётки и измерение электросопротивления.
\end{orientir}

\subsubsection{Прямое наблюдение}

Прибор, позволяющий видеть отдельные атомы и вакансии,~—
\term{автоионный (полевой ионный) микроскоп}; его разрешение
$0{,}2$--$0{,}3$~нм. Однако у метода серьёзные ограничения:
\begin{itemize}
\item видны только атомы на поверхности острия, а не в объёме;
\item метод применим лишь к металлам, у которых поле, необходимое для
отображения, не вызывает их собственного испарения (тугоплавкие металлы);
\item сильное электрическое поле создаёт в образце большие механические
напряжения, поэтому истинные искажения вокруг дефекта установить нельзя.
\end{itemize}

\subsubsection{Косвенные методы: что чувствует дефекты}

Многие свойства зависят от концентрации точечных дефектов:
\begin{itemize}
\item дефекты и искажённые области вокруг них рассеивают электроны и
нейтроны~— меняются электросопротивление и прохождение нейтронов;
\item меняются параметр решётки и плотность;
\item увеличивается запасённая энергия (у каждого дефекта есть энергия
образования);
\item меняются упругие модули и, через взаимодействие с дислокациями,
механические свойства.
\end{itemize}
Измеряя такое свойство при разных температурах или в ходе нагрева,
находят энергии образования и миграции дефектов.

\subsubsection{Метод Симмонса--Баллуффи: длина образца и параметр решётки}

\paragraph{Идея}
При образовании вакансии атом из объёма переносится на поверхность.
Число узлов растёт на единицу, и объём кристалла увеличивается на один
атомный объём. Если бы атомы вокруг вакансии не смещались, межатомные
расстояния (и параметр решётки) не изменились бы, а длина образца
выросла бы. Поэтому \emph{разница} между относительным удлинением
образца и относительным изменением параметра решётки и есть мера числа
вакансий.

Релаксация вокруг дефекта немного меняет параметр решётки. Эшелби
показал, рассматривая кристалл как изотропную упругую среду, что
искажение вокруг точечного дефекта (центра давления) меняет параметр
решётки ровно так, как это соответствует изменению объёма от релаксации.
Если $V_x$~— изменение объёма из-за релаксации вокруг одного дефекта,
$C$~— атомная доля дефектов, $V_0$~— объём на атом, то
\[
\frac{\Delta a}{a}=\frac13\cdot\frac{V_x}{V_0}\,C .
\]

\begin{shagi}
\item Объём на один атом кристалла с дефектами:
$V=V_0(1\pm C)+V_xC$ (знак «$+$» для вакансий, «$-$» для межузельных
атомов: каждая вакансия добавляет узел, каждый межузельный атом его
«отнимает»).
\item Относительное изменение объёма:
$\dfrac{\Delta V}{V_0}=\pm C+\dfrac{V_x}{V_0}C=\pm C+3\dfrac{\Delta a}{a}$.
\item Для малых деформаций $\Delta V/V_0=3\Delta L/L$, где $L$~— длина
образца.
\end{shagi}
Отсюда
\begin{keyeq}\label{eq:simmons}
C=\pm3\left(\frac{\Delta L}{L}-\frac{\Delta a}{a}\right).
\end{keyeq}

\begin{figure}[ht]
\centering
\begin{tikzpicture}
\begin{axis}[width=0.66\textwidth, height=5.6cm,
  xmin=600, xmax=940, ymin=1.0, ymax=2.6,
  xlabel={$T$, К}, ylabel={относительное изменение, \%},
  legend pos=north west, legend cell align=left]
  \addplot[fmRed, very thick, domain=600:933, samples=60]
    {1.0+0.0028*(x-600)+0.0000035*(x-600)^2 + 8*100*exp(1.0)*exp(-0.67/(8.617e-5*x))/3};
  \addlegendentry{$\Delta L/L$ (длина образца)}
  \addplot[fmBlue, very thick, domain=600:933, samples=60]
    {1.0+0.0028*(x-600)+0.0000035*(x-600)^2};
  \addlegendentry{$\Delta a/a$ (рентген)}
  \draw[dimline] (axis cs:928,2.29) -- (axis cs:928,2.46);
  \node[font=\footnotesize, anchor=east] at (axis cs:926,2.40) {$C_v/3$};
\end{axis}
\end{tikzpicture}
\caption{Схема опыта Симмонса--Баллуффи (алюминий). При низких температурах
длина образца и параметр решётки растут одинаково (тепловое расширение);
у точки плавления длина растёт быстрее: в кристалле появляются вакансии.
Разность кривых равна $C_v/3$ (на рисунке она увеличена в 8 раз; у точки
плавления она составляет лишь $\approx0{,}02\,\%$ при тепловом расширении
$\approx1{,}5\,\%$)}
\label{fig:simmons}
\end{figure}

Опыт ставят так: одновременно измеряют удлинение образца (дилатометром)
и параметр решётки (рентгенографически) при нагреве до температуры,
близкой к плавлению (рис.~\ref{fig:simmons}). Метод хорош тем, что даёт
\emph{абсолютную} концентрацию вакансий, не требуя знания других
характеристик. Именно так впервые было надёжно доказано, что
равновесными дефектами в металлах являются вакансии (для Al у точки
плавления $C_v\approx9\cdot10^{-4}$).

\subsubsection{Электросопротивление}

Электросопротивление металла складывается из тепловой части (рассеяние
электронов на колебаниях решётки) и \term{остаточного сопротивления},
которое создают примеси и дефекты (подробнее~— в разделе~3.5). Любые
искажения решётки рассеивают электроны проводимости, поэтому остаточное
сопротивление очень чувствительно к дефектам. При малой концентрации
добавка к сопротивлению \emph{пропорциональна} концентрации дефектов:
\[
\Delta\rho=\rho_1C,
\]
где $\rho_1$~— вклад одного процента (или одной доли) дефектов; он
зависит от типа дефекта и металла. Изменение остаточного сопротивления~—
удобная мера концентрации дефектов. Вклад дефектов в тепловую часть мал
и плохо известен, поэтому измерения обычно ведут при низкой температуре
(например, в жидком гелии).

\paragraph{Измерение при высокой температуре}
При нагреве сопротивление растёт почти линейно. Однако при высоких
температурах экспериментальные точки отклоняются от прямой вверх
(рис.~\ref{fig:rho-T},\,а): добавка $\Delta\rho$ создаётся равновесными
вакансиями. Поскольку $C_v\propto\exp(-\EvF/kT)$,
\[
\Delta\rho=A\Ex{-\frac{\EvF}{kT}},\qquad \ln\Delta\rho=\text{const}-\frac{\EvF}{k}\cdot\frac1T .
\]
В координатах $\ln\Delta\rho$\,--\,$1/T$ (координатах Аррениуса) получается
прямая; по её наклону находят $\EvF$ (рис.~\ref{fig:rho-T},\,б), а затем
и саму концентрацию.

\begin{figure}[ht]
\centering
\begin{subfigure}[b]{0.48\textwidth}\centering
\begin{tikzpicture}
\begin{axis}[width=\linewidth, height=5cm, xmin=0, xmax=1, ymin=0, ymax=1.25,
  xlabel={$T$}, ylabel={$\rho$}, xtick=\empty, ytick=\empty]
  \addplot[fmBlue, dashed, thick, domain=0.1:1] {0.05+0.9*x};
  \addplot[fmRed, very thick, domain=0.1:1, samples=80] {0.05+0.9*x+0.25*exp(-10*(1-x))};
  \draw[dimline] (axis cs:0.97,0.923) -- (axis cs:0.97,1.108);
  \node[font=\small, anchor=east] at (axis cs:0.95,1.03) {$\Delta\rho$};
\end{axis}
\end{tikzpicture}
\caption{}
\end{subfigure}\hfill
\begin{subfigure}[b]{0.48\textwidth}\centering
\begin{tikzpicture}
\begin{axis}[width=\linewidth, height=5cm, xmin=0, xmax=1, ymin=0, ymax=1,
  xlabel={$1/T$ (или $1/T_{\text{з}}$)}, ylabel={$\ln\Delta\rho$}, xtick=\empty, ytick=\empty]
  \addplot[black, thick, domain=0.05:0.95] {0.9-0.8*x};
  \addplot[only marks, mark=*, mark size=1.8pt, fmRed] coordinates {(0.15,0.79) (0.3,0.66) (0.48,0.51) (0.62,0.41) (0.8,0.25)};
  \draw (axis cs:0.4,0.58) -- (axis cs:0.65,0.58) -- (axis cs:0.65,0.38);
  \node[font=\footnotesize, anchor=west] at (axis cs:0.66,0.5) {наклон $-\EvF/k$};
\end{axis}
\end{tikzpicture}
\caption{}
\end{subfigure}
\caption{Определение энергии образования вакансий по электросопротивлению:
\textit{а}~— отклонение сопротивления от линейной зависимости при высоких
температурах; \textit{б}~— та же добавка в координатах Аррениуса}
\label{fig:rho-T}
\end{figure}

\paragraph{Закалочные опыты}
Образцы закаливают с разных температур $T_{\text{з}}$ и измеряют прирост
остаточного сопротивления $\Delta\rho$ при низкой температуре. Избыточная
концентрация вакансий равна равновесной при $T_{\text{з}}$, поэтому
\[
\ln\Delta\rho=\text{const}-\frac{\EvF}{kT_{\text{з}}} ,
\]
и снова наклон прямой в координатах $\ln\Delta\rho$\,--\,$1/T_{\text{з}}$ даёт
$\EvF$. Затем образцы отжигают и по скорости убывания $\Delta\rho$
методами п.~\ref{ssec:annealing} находят энергию миграции $\Evm$.
Сумма $\EvF+\Evm$ должна совпасть с энергией активации самодиффузии,
измеренной независимо,~— это важная проверка всей картины.

\begin{table}[ht]
\centering
\caption{Сравнение высокотемпературных и закалочных измерений}
\label{tab:exp-compare}
\small
\begin{tabularx}{\textwidth}{@{}l>{\raggedright\arraybackslash}X>{\raggedright\arraybackslash}X@{}}
\toprule
 & Достоинства & Недостатки \\
\midrule
Закалка & высокая точность измерений при низкой температуре &
  нельзя измерить исходную концентрацию моновакансий: при закалке часть
  вакансий уходит на стоки, образует дивакансии и комплексы \\
Высокотемпературные измерения & измеряется истинно равновесное состояние &
  ниже точность: нужно выделить малую добавку на фоне большого теплового
  вклада \\
\bottomrule
\end{tabularx}
\end{table}

\begin{podrobnee}[Современные методы]
Сегодня важнейший метод исследования вакансий~— \emph{позитронная
аннигиляция}: позитрон захватывается вакансией (там нет положительного
ядра), и его время жизни и спектр аннигиляционного излучения меняются.
Метод очень чувствителен к вакансиям ($C_v\gtrsim10^{-6}$) и позволяет
различать моновакансии и их скопления. Широко используется и
просвечивающая электронная микроскопия (видны скопления дефектов~—
петли, поры), а также компьютерное моделирование из первых принципов.
\end{podrobnee}

\begin{itogi}
\begin{itemize}
\item Прямое наблюдение (автоионный микроскоп) ограничено поверхностью и
сильными полями; основная информация~— из косвенных методов.
\item Метод Симмонса--Баллуффи: $C=\pm3(\Delta L/L-\Delta a/a)$~—
абсолютная концентрация вакансий.
\item Остаточное сопротивление пропорционально концентрации дефектов;
в координатах $\ln\Delta\rho$\,--\,$1/T$ наклон даёт $\EvF$.
\item Отжиг закалённых образцов даёт $\Evm$; проверка: $\EvF+\Evm=\Ecd$.
\end{itemize}
\end{itogi}

\begin{proverka}
\item Почему при образовании вакансий длина образца растёт быстрее, чем
параметр решётки?
\item Какой знак имеет разность $\Delta L/L-\Delta a/a$, если в кристалле
преобладают межузельные атомы?
\item Почему остаточное сопротивление измеряют при низких температурах?
\item Почему в закалочных опытах нельзя точно измерить исходную
концентрацию моновакансий?
\end{proverka}
